\section{A fourth-order growth estimate}\label{sec:growth}

The reach and interface estimates are local.  We now integrate them to
bound the portions of the zero set on which the radial argument will
incur errors: large defect, appreciable curvature, or short reach.
Throughout this section $v:\R^7\to(-1,1)$ is an entire stable solution
of $\Delta v=v^3-v$, $D>0$, and its energy density is unbounded, as in
Sections~\ref{sec:reach} and~\ref{sec:interface}.  In particular $m=6$.
Recall that $a=|A|$, $d$ is the penalized reach, and $G_1$ is the
constant value of its angular penalty near normal incidence.  Write
\[
 K(R)=1+\sup_{x\in\R^7}M_R(x),\qquad
 \kappa=\frac{\sqrt2}{G_1},\qquad E=e^{-\kappa d}.
\]
The target scale is $K(R)R^{m-4}$ for the fourth curvature moment and
the squared defect and interaction.  Its gain of four powers over the
weighted area scale $K(R)R^m$ will pay for rough regions and curvature
errors in Section~\ref{sec:radial}, even when $K(R)$ grows arbitrarily
slowly.  There are three steps.  Simons' identity reduces the curvature
moment to a small multiple of the squared defect.  A single ambient
stability test then supplies both surface mass and mass near bad zeros.
Finally, a multiplier for the reach inequality bounds the derivative
cost of that test.  The integral comparison of
Lemma~\ref{lem:interface-integral} will close the three estimates.

All surface integrals in this section use induced area.  Fix a threshold
$\delta\in(0,1/4)$ provisionally and define
\[
 X=\{z\in\Sigma:D(z)\ge\delta\},\qquad
 \Sigma_g=\Sigma\setminus X,
\]
and let $\mathcal X\subset X$ be any maximal unit-separated set.  Put
\[
 \chi(x)=\chi_{R,x_c}(x)
 =\exp\!\left(-\frac{\sqrt{1+|x-x_c|^2}}R\right),\qquad
 N_R=\sum_{x\in\mathcal X}\chi(x)^2,
 \qquad S_0=\frac{a^2+I_P}{2}.
\]
Here $I_P$ is the normal average defined in
Lemma~\ref{lem:interface-sharp-means}; its finite depth $P$ will be fixed
below.  Thus $S_0$ is the favorable mass in the strengthened surface
inequality.  The convention when $X=\varnothing$ is $N_R=0$ and
$\dist(x,X)=+\infty$.  Since $D=1$ at a singular zero, every such zero
belongs to $X$.  The quantity $N_R$ records the bad set through any
maximal unit-separated net, rather than requiring that set to be smooth.

\begin{proposition}[Fourth-order growth]\label{prop:growth-fourth-order}
In the seven-dimensional stable setting above, fix the penalty $\Gamma$ sufficiently close to the choices specified in
Section~\ref{sec:reach}, with $\eta_0<10^{-4}$ and
$G_1-1,(1-\delta_2)^{-1}-1<10^{-3}$, decreasing these errors if required
by Section~\ref{sec:radial}.  There are a finite $P$, a positive
$\delta$, and constants $R_0,C<\infty$ such that, for every solution in
the setting above, every $x_c\in\R^7$, every maximal net
$\mathcal X$, and every $R\ge R_0$,
\begin{equation}\label{eq:M1}
 \mathcal Y_R:=
 \int_{\Sigma_g}\chi^2(a^4+D^2+E^2)+N_R
 \le C K(R)R^{m-4}.
\end{equation}
The constants depend on the fixed penalty and tolerances, including
$\delta$, but are independent of the solution, the center, the net,
and $R$.  Any prescribed smaller positive tolerances in
\eqref{eq:L1}--\eqref{eq:L4} and \eqref{eq:L12}--\eqref{eq:L13}
may be imposed before $P$ and $\delta$ are fixed.
\end{proposition}

\paragraph{Order of choices.}
The threshold $\delta$ must remain fixed while the spatial radius grows;
otherwise the bad set and its net could change with $R$.  First fix the penalty
and the threshold-independent constant $C_D$ in \eqref{eq:L11}.
Choose the small coefficient in the curvature estimate below in terms
of $C_D$.  Next fix the interface tolerances, including
$\tau\le .001$ for Lemma~\ref{lem:growth-gluing}, the normal depth $P$,
and the required lower reach threshold $t_{\min}$.  Only then shrink
$\delta$ so that the interface estimates hold on $\Sigma_g$ and
$d\ge t_{\min}$ there.  In particular, throughout the remaining proof,
\begin{equation}\label{eq:growth-sharp-budget}
 I:=I_P\ge .999a^2,\qquad c_I D\le I\le C_P D,
 \qquad |H|\le\left(\frac1{2\sqrt2}+\eps_H\right)I,
 \quad \eps_H\le .001.
\end{equation}
The value of $\eps_H$ can also be small enough that $|H|\le .355I$.
After $\delta$ is fixed, choose a collar width $L_0$ about $X$.
The reach is bounded on that collar, so a compact interval
$[t_{\min},t_*]$ can cover all its reach values.  The later test amplitude
$f(d)$ will equal one on this interval.  Its decay and the parameters
that join its different decay ranges are chosen last.  All these choices
precede every spatial exhaustion and the limit $R\to\infty$.

\paragraph{Area and collar bounds.}
Two elementary estimates will justify all weighted integrations:
\begin{equation}\label{eq:growth-area}
 K(t)\le C\frac tR K(R)\quad(t\ge R\ge1),\qquad
 \int_{\Sigma_g}\chi^2\le CK(R)R^m.
\end{equation}
To prove the first, cover a ball of radius $t$ by at most
$C(t/R)^n$ balls of radius $R$ and divide the resulting energy bound by
$t^m$.  For the second, the uniformly regular small charts at good zeros
bound their area by the energy in fixed enlargements; use a bounded
overlap covering and then exponentially weighted annuli.  The same
charts, or Lemma~\ref{lem:interface-good-area}, imply
\begin{equation}\label{eq:growth-collar-count}
 \int_{\Sigma_g\cap\{\dist(\cdot,X)\le s\}}\chi^2
 +\int_{\{\dist(\cdot,X)\le s\}}\chi^2\,\dd x
 \le C(1+s)^n e^{C(1+s)/R}N_R\quad(s\ge1).
\end{equation}
Indeed maximality of the net places $X$ in its unit balls, and each
enlarged ball has good area and volume at most $C(1+s)^n$.  Shifting the
weight through that ball costs at most $e^{C(1+s)/R}$.  These arguments
also prove finiteness of every integral used below.

\subsection{Simons' identity and the curvature moment}

We first estimate $a^4$ with a coefficient of $D^2$ that can be made
arbitrarily small.  The method combines Simons' identity, a strict Kato
gain, and stability, following the curvature-estimate arguments of
\citet{Simons1968,SchoenSimonYau1975}.  Our surface need not be minimal.
We therefore retain both the derivatives of $H=\operatorname{tr}A$ and
the cubic-curvature error involving $H$; the averaged mean-derivative
estimate \eqref{eq:L2} will make their cost sufficiently small.

\begin{lemma}[Curvature controlled by a small defect cost]
\label{lem:growth-curvature}
For each $\eps_2>0$ there are choices of the interface tolerances and
$\delta>0$ such that, uniformly for all centers and sufficiently large
$R$,
\begin{equation}\label{eq:M2}
 A_4:=\int_{\Sigma_g}\chi^2a^4
 \le\eps_2\int_{\Sigma_g}\chi^2D^2
       +C_{\eps_2,\delta}N_R
       +C_{\eps_2,\delta}K(R)R^{m-4}.
\end{equation}
The constants remain valid on decreasing $\delta$ once the preceding
interface tolerances have been fixed, except for the displayed allowed
dependence of the last two coefficients on $\delta$.
\end{lemma}

\begin{proof}
The strict gradient gain lets Simons' identity and surface stability
control one another.  We derive that gain first, apply it away from
$X$, and then charge the omitted collar to $N_R$.

\paragraph{The Kato gain.}
We first record the algebra behind the strict gradient gain.  If $A$ is
a symmetric tensor and $\nabla A$ is totally symmetric, there are
dimensional $\theta>0,C$ such that
\begin{equation}\label{eq:growth-kato}
 |\nabla A|^2\ge(1+\theta)|\nabla|A||^2-C|\nabla H|^2.
\end{equation}
For $|A|=|\nabla A|=1$, equality in ordinary Kato requires
$\nabla_iA=b_iA$.  Total symmetry, after placing $b$ in the first
coordinate direction, forces $A$ to have only its $11$ component.
Consequently $\nabla H=b\tr A\ne0$.  Compactness of the normalized
finite-dimensional tensor sets therefore gives a strict Kato gap when
$|\nabla H|$ is sufficiently small; when it is larger, the last term in
\eqref{eq:growth-kato} pays for the inequality.  Homogeneity proves the
claim.  The weak gradient of $|A|$ is zero almost everywhere on its zero
set, so this includes that set.

\paragraph{Simons' identity away from bad zeros.}
Choose the averaging radius $S$ in \eqref{eq:L2}.  Let $m_c$ be a
Lipschitz cutoff vanishing within distance $3S$ of $X$, equal to one
outside distance $6S$, and with bounded gradient; take $m_c=1$ if $X$ is
empty.  Put $u=\chi m_c$.  The hypersurface identity, with our convention
$A=\dd\nu$, is
\[
 \Delta_\Sigma A=\nabla_\Sigma^2H+HA^2-a^2A.
\]
Pair with $A u^2$ and integrate by parts.  Codazzi gives
$\operatorname{div}_\Sigma A=\nabla_\Sigma H$.  Thus the term with
$\nabla^2H$ costs at most
$C\int u^2|\nabla H|^2+C\int a^2|\nabla u|^2$; the cross from
$\Delta A$ uses $\langle A,\nabla_iA\rangle=a\nabla_i a$ and costs an
arbitrarily small fraction of $\int u^2|\nabla a|^2$ plus
$C\int a^2|\nabla u|^2$.  Since $|H|\le CD$, \eqref{eq:growth-kato}
yields, for a dimensional $\theta'>0$,
\begin{equation}\label{eq:growth-simons}
 (1+\theta')G\le F+
 C\int u^2\bigl(|\nabla H|^2+Da^3\bigr)
 +C\int a^2|\nabla u|^2,
 \quad G=\int u^2|\nabla a|^2,\quad F=\int u^2a^4.
\end{equation}
Applying \eqref{eq:L3} to $au$, and using the curvature lower bound in
\eqref{eq:L1} with sufficiently small tolerance, gives
\[
 F\le(1+\theta'/2)G+C\int a^2|\nabla u|^2.
\]
Subtracting this from \eqref{eq:growth-simons} controls both $G$ and $F$
by the remaining errors.

\paragraph{Paying the mean-curvature and collar errors.}
It remains to control the errors in \eqref{eq:growth-simons} without
losing the small coefficient of $D^2$.  For every center in this
application there is one uniform sheet on $B_{4S}(z)$.  Fubini in
\eqref{eq:L2} therefore gives
\[
 \int u^2|\nabla H|^2
 \le C\eps\int_{\Sigma_g}\chi^2D^2.
\]
The averaging balls do not meet $X$.  For each fixed point counted by
the integral, its centers lie on one chart of area at most $CS^m$;
the factor $S^{-m}$ cancels this bound.  The ratio of the two weights is
at most $e^{2S/R}$, hence at most two for sufficiently large $R$.  This
explains why the constant multiplying $\eps$ is independent of $S$.
Furthermore $a^2\le CD$ implies $Da^3\le C\sqrt\delta D^2$.
The cutoff cost is bounded by
\[
 C_S N_R+CR^{-2}\int_{\Sigma_g}\chi^2a^2
 \le C_SN_R+\eps' A_4
      +C_{\eps'}R^{-4}\int_{\Sigma_g}\chi^2.
\]
The discarded part of $A_4$, where $m_c\ne1$, is also $C_SN_R$ by
\eqref{eq:growth-collar-count}.  Choose $\eps'$ small enough to absorb,
then the tolerance in \eqref{eq:L2} and $\delta$ small enough for the
prescribed $\eps_2$.  Equation~\eqref{eq:growth-area} proves
\eqref{eq:M2}.  All integrations may first have compact support; the
exponential weight, polynomial area bound, and uniform unit derivative
bounds remove the additional cutoffs.  Approximation of $a$ by
$(a^2+\varepsilon^2)^{1/2}$ justifies its use as a test.
\end{proof}

From now on fix $\eps_2$ sufficiently small relative to $C_D$ in
\eqref{eq:L11}, and make all the preceding parameter choices.  Lower
bounds on $t_{\min}$ used below involve only the already fixed reach and
interface constants.  Lemma~\ref{lem:reach-small-defect} permits us to
enforce them by the final choice of $\delta$.

\subsection{Gluing ambient and surface stability}

The curvature estimate leaves a cost proportional to $N_R$.  To
control it, stability must detect defect near $X$, where a regular
surface test is unavailable.  The ambient amplitude $q=1-v^2$ does
this.  Away from $X$, lifting the normal profile
$G(s)=1-g(s)^2$ supplies the stronger surface mass $S_0$.
We now combine those amplitudes into one physical test.

Combining ambient and transition-set tests is also part of the
free-boundary methods of \citet{CFFS2025} and the smooth quartic
three-dimensional argument of \citet[Sections~3 and~5]{LLWWW2026}.
The latter uses a sum of transition-profile derivatives and a linear
joining to the ambient amplitude.  A related geometric stability
inequality in \citet[Proposition~4.7]{CFFFS2026} uses a weighted
ground-state identity on the normal parametrizing manifold.  Here we
use a square root of a sum of squares and retain its exact gradient
variance.  The following
argument supplies the required estimates on our six-dimensional
interface.

The cancellation is already visible for two nonnegative Lipschitz
functions $\xi_0,\xi_1$ on a Euclidean open set.  Wherever
$\xi_0^2+\xi_1^2>0$, differentiation gives
\[
 |\nabla\xi_0|^2+|\nabla\xi_1|^2
 -\left|\nabla\sqrt{\xi_0^2+\xi_1^2}\right|^2
 =\frac{|\xi_0\nabla\xi_1-\xi_1\nabla\xi_0|^2}
        {\xi_0^2+\xi_1^2}.
\]
The identity holds almost everywhere with value zero on their common
zero set.  If $\xi_0=h\cos\theta$ and $\xi_1=h\sin\theta$ for
locally Lipschitz $h\ge0$ and $0\le\theta\le\pi/2$, the square root is just $h$ and
the variance is $h^2|\nabla\theta|^2$.  Thus it cancels the cost of
changing amplitudes when the profiles agree.  In the construction
below, $q$ and the lifted $G$ agree to a controlled error on the short
normal segments called accurate.  The other short segments are charged
to defect mass, while the long segments have exponentially small tails.
This is a change from an ambient description to a surface description
near bad zeros.  The interpolation between neighboring sheet amplitudes
in Section~\ref{sec:radial} is a separate construction.

\Needspace{6\baselineskip}
\begin{lemma}[Uniform gluing]\label{lem:growth-gluing}
Fix $0<\tau\le .001$ among the interface tolerances, before the final
choice of $\delta$, and choose the preceding data with this tolerance.
There exist
$L_0,t_*,c_0>0$ and $R_1,C<\infty$ with the following property.
Let $f:[t_{\min},\infty)\to(0,1]$ be nonincreasing and locally Lipschitz,
equal to one on $[t_{\min},t_*]$, and a positive multiple of
$e^{-\kappa t/2}$ for sufficiently large $t$.  Suppose also that
$|f'|\le C_f^{(0)}f$ for some finite constant.  Then for every center
and $R\ge R_1$, writing $s_X=\dist(\cdot,X)$,
\begin{equation}\label{eq:M3}
 \begin{aligned}
 &(1-\tau)\int_{\Sigma_g\cap\{s_X\ge4L_0\}}S_0 f(d)^2\chi^2+c_0N_R\\
 &\qquad\le(1+\tau)\int_{\Sigma_g}\chi^2 f'(d)^2|\nabla_\Sigma d|^2
       +\frac C{R^2}\int_{\Sigma_g}\chi^2f(d)^2.
 \end{aligned}
\end{equation}
In particular $c_0$ is independent of $f$.  The constants $L_0,t_*,R_1,C$
can also be chosen independently of $f$; its stated tail and slope
conditions are used only to justify exhaustion.
\end{lemma}

\begin{proof}
If $X$ is empty, apply \eqref{eq:L3} to $\chi f(d)$ and use Young's
inequality.  Suppose henceforth that $X$ is nonempty.  We construct the
test, reserve positive mass for its transition region, and estimate its
derivatives.  The $f'$ energy and the spatial-weight cost remain on the
right side of \eqref{eq:M3}.  After estimating those terms and using the
mass reserve, only the two partition terms remain to be canceled by
the variance identity.

\paragraph{Construction: a collar on which $f=1$.}
The base tangent-ball estimate in Corollary~\ref{cor:reach-base-defect}
and \eqref{eq:diffuse-local} give a uniform bound
\begin{equation}\label{eq:growth-bounded-collar-reach}
 \sup_{\Sigma_g\cap\{s_X\le10L_0\}}d\le T(\delta,L_0)<\infty.
\end{equation}
Indeed a nearby $x\in X$ gives
$D(z)\ge c\delta e^{-C(10L_0+1)}$, whereas a base tangent ball gives
$D(z)\le C_\Gamma/d(z)$ when $d$ is large.  The same bound, with a
changed constant, applies to the \emph{source} reach $d(z)$ of a
plateau pair $(z,y)$ whenever $s_X(y)\le10L_0$, because
\eqref{eq:C10} gives $D(y)\le C_\Gamma/d(z)$.  Choose $t_*$ to cover
both bounds.  This choice will also be used in the multiplier argument.

Choose Lipschitz functions on physical space
\[
 \begin{gathered}
 \theta=\theta(s_X/L_0),\qquad
 \theta=0\ (s_X\le2L_0),\qquad
 \theta=\pi/2\ (s_X\ge3L_0),\\
 |\nabla\theta|\le C/L_0,\qquad
 \alpha_0=\cos\theta,\qquad b=\sin\theta.
 \end{gathered}
\]
On $\Sigma_g$ choose $\mu=0$ for $s_X\le L_0/2$ and $\mu=1$ for
$s_X\ge L_0$, with $|\nabla\mu|\le C/L_0$, and put
$p_0(z)=\mu(z)f(d(z))$.  Use the normal immersion
$\Phi(z,s)=z+s\nu(z)$ and the even curvature cutoff $\beta(z,s)$ from
Lemma~\ref{lem:interface-immersed-test}, supported strictly inside
$|s|a<\alpha$ and equal to one for $|s|a\le\alpha/2$.  The parameter
$\alpha$ is fixed small in terms of $\tau$ before shrinking $\delta$.
The physical ambient function and the lifted functions are
\[
 \xi_0(x)=q(x)\chi(x)\alpha_0(x),\qquad
 \xi_{z,s}(x)=G(s)\beta(z,s)\chi(x)b(x)p_0(z),\quad x=\Phi(z,s).
\]
They define the single physical test $F=(\sum_j\xi_j^2)^{1/2}$.
Here the sum includes $\xi_0$ and every nonzero lifted preimage of the
physical point; no global numbering of the sheets is chosen.

\paragraph{Admissibility and the full variance.}
The immersion need not be globally injective.  To justify the physical
test without choosing global sheet labels, first cap
$|s|$ and physical support by smooth cutoffs.  The nonzero lifted support
is then compact inside the immersion domain: $|z|$ is bounded because
both $|x|$ and $|s|$ are bounded, and $\mu$ keeps the bases away from
$X$.  Fibers are discrete by local invertibility and finite by
compactness.  Around any physical point finitely many local inverses
therefore contain all nonzero branches.  A smaller neighborhood has no
additional branch arriving from elsewhere in the compact support.
Thus $F$ is locally Lipschitz, including where it vanishes, and is a
compactly supported stability test.  Almost everywhere,
\begin{equation}\label{eq:growth-variance}
 \sum_j|\nabla\xi_j|^2-|\nabla F|^2
 =\min_{w\in\R^n}\sum_j|\nabla\xi_j-\xi_jw|^2
 =:\mathcal V\ge0.
\end{equation}
The potential terms add exactly, even where $3v^2-1$ is negative.
Consequently the physical quadratic form is the sum of the ambient and
lifted forms \emph{minus} $\int\mathcal V$.  No partition of a chart is
introduced.  The area formula accounts for every lifted preimage.
Moreover the variance of any subfamily is a lower bound for
$\mathcal V$, since each omitted square is nonnegative before
minimization.

For each fixed normal cap remove the spatial cap, and then remove the
normal cap.  For the lifted terms this is dominated convergence in
\eqref{eq:L4}, using
\[
 \chi(\Phi(z,s))/\chi(z)\le e^{|s|/R},\qquad
 |\nabla d|\le C\exp\bigl(C\sqrt{\log(K(d)+1)}\bigr),\qquad
 K(d)\le1+Cd.
\]
For each fixed $f$ the products $|f'|\,|\nabla d|$ are bounded, and normal cutoff costs are
polynomials in $|s|$ times $G(s)^2$.  Polynomial good-area growth and
the exponential spatial weight dominate the remaining factors.
For the subtracted variance retain only the lower bounds on fixed
short segments used below, and apply Fatou's lemma.  This proves that
the same stability comparison is valid after both caps are removed.
The dominating bound here may depend on $f$; it is used only for
convergence.  No coefficient in the quantitative estimates below uses
its value.

\paragraph{Claim 1: a positive mass reserve independent of $f$.}
The favorable terms will pay for the transition collar while retaining
the full stated strength outside it.  Define
\[
 S_{\rm sw}=\Sigma_g\cap\{3L_0/2\le s_X\le4L_0\},\qquad
 D_{\rm sw}=\int_{S_{\rm sw}}\chi^2D.
\]
Disjoint fixed small balls about the unit-separated bad net yield
$c_\delta N_R$ from the term
$\int q^3D\chi^2\alpha_0^2$: at their centers $q=1$ and $D\ge\delta$,
and unit regularity and \eqref{eq:diffuse-local} give uniform lower
bounds in the balls.  These balls are disjoint from a fixed small
normal tube over $S_{\rm sw}$.  That tube is injective: small defect
gives a single graph in a uniform fixed neighborhood of every base,
and any overlapping short normal segments must be in such a chart,
where the normal map is injective.  In this tube the ambient term gives
\[
 c\int_{S_{\rm sw}}\chi(z)^2\alpha_0(z)^2D(z)
       -\frac C{L_0}D_{\rm sw}-o_{R\to\infty}(1)D_{\rm sw}.
\]
Here $q^3D$ is uniformly comparable to $D(z)$, and freezing $\alpha_0$
and $\chi$ accounts for the two errors.

The lifted potential and Jacobian terms of \eqref{eq:L4} provide the
complementary positive mass with $b(z)^2$, and on $s_X\ge4L_0$ provide
$(1-\tau')\sigma\int S_0f^2\chi^2$, where $\tau'>0$ is as small as
prescribed.  To check the freezing error, on the cutoff tube
\[
 |\partial_sJ|\le C(1+|s|)^nD(z),\qquad J=\det(1+sA),
\]
because $|H|\le CD$ and $a^2\le CD$; expand the determinant in powers
of $sA$.  The nonpositive potential term can be restricted to the
fixed finite segments used to extract $I_P$.  Freezing $b^2,\chi^2$ on
these terms therefore costs at most
$C L_0^{-1}D_{\rm sw}+o_R(1)D_{\rm sw}$ on $S_{\rm sw}$.
Outside $S_{\rm sw}$ any change of $b$ between a supported base and its
image requires
\[
 |s|\ge L_0/2,\qquad s_X(z)\le3L_0+|s|.
\]
Equation~\eqref{eq:growth-collar-count} and exponential decay of $G$
bound the resulting integrals by $o_{L_0}(1)N_R$, uniformly for large
$R$.  The curvature cutoff errors $Ce^{-c/a}f^2$ are absorbed by the
available $I f^2$ on shrinking $\delta$.  Here is an explicit mass
reserve.  Let $c_{\rm amb}>0$ denote the ambient switch coefficient,
take the pure-sheet loss below $\tau/16$, and fix
\[
 0<c_{\rm sw}\le
 \min\{c_{\rm amb}/4,\,\tau\sigma c_I/32\}.
\]
On $S_{\rm sw}\cap\{s_X<4L_0\}$ the complementary ambient and sheet
gains supply at least $2c_{\rm sw}(\alpha_0^2+b^2)D=2c_{\rm sw}D$.
On $s_X=4L_0$, retain $(1-\tau/2)\sigma S_0f^2$ for exterior strength.
Since $b=f=1$, the spare sheet gain is at least
$(7\tau/16)\sigma S_0\ge(7\tau\sigma c_I/32)D$, which also pays
$2c_{\rm sw}D$.  This leaves $c_{\rm sw}D_{\rm sw}$ for errors and
the desired exterior strength after the remaining tolerance losses.
This treatment includes the level set $s_X=4L_0$, even if it has
positive surface area.  We have therefore reserved positive fixed
multiples of both $D_{\rm sw}$ and $N_R$; their coefficients were
chosen before $f$.  The next two claims show that this reserve suffices.

\paragraph{Claim 2: every derivative cost except the partition terms.}
For clarity we expand every differentiated amplitude in the lifted
gradient term.  Let $B_s=(\Id+sA)^{-1}$ on tangent vectors and write
$\beta=\beta_0(sa)$ with $\beta_0$ even.  Almost everywhere,
\[
 \nabla\beta=\beta_0'(sa)
       \bigl(a\nu+sB_s\nabla_\Sigma a\bigr),
 \qquad |\nabla\beta|\le C_\alpha(1+|s|).
\]
This includes the derivative of $a$; its weak gradient is uniformly
bounded by the unit derivative estimates on good charts.  At $a=0$
the derivative of this cutoff is zero almost everywhere.
All vectors in the following expansion are physical vectors at
$x=\Phi(z,s)$:
\[
 \begin{aligned}
 V_f&=\chi b\beta\mu f'(d)B_s\nabla_\Sigma d,&
 V_\mu&=\chi b\beta f B_s\nabla_\Sigma\mu,\\
 V_\beta&=\chi b\mu f\nabla\beta,&
 V_\chi&=b\beta\mu f\nabla\chi,\\
 V_b&=\chi\beta\mu f\nabla b.&
 \end{aligned}
\]
Their sum is $\nabla(\chi b\beta\mu f)$.
Keep $|V_b|^2+2V_b\cdot V_\chi$ untouched, together with the ambient
partition terms.  These are the following ambient and lifted
densities, respectively:
\begin{equation}\label{eq:M4}
 \begin{split}
 &\chi^2q^2|\nabla\alpha_0|^2
      +2\chi q^2\alpha_0\nabla\chi\cdot\nabla\alpha_0,\\
 &G^2Jp_0^2\beta^2
   \bigl(\chi^2|\nabla b|^2+2\chi b\nabla\chi\cdot\nabla b\bigr).
 \end{split}
\end{equation}
For a fixed arbitrarily small $\varepsilon>0$, all other squares and
crosses are bounded by
\[
 (1+C\varepsilon)|V_f|^2
 +C_\varepsilon\bigl(|V_\mu|^2+|V_\beta|^2+|V_\chi|^2\bigr)
 +2V_b\cdot(V_f+V_\mu+V_\beta).
\]
The vector $V_b$ differentiates the physical partition.  Its displayed
square and its cross with $V_\chi$ must keep their signs for the later
cancellation; all other terms can be estimated now.  Every lifted
integral below has the factor $G^2J$.
The $V_f$ square uses the inverse metric bound $1+O(\alpha)$ and the
normal integral $\sigma+O(R^{-1})$.  The $V_\chi$ square uses
$|\nabla\chi|\le\chi/R$.  The $V_\mu$ square is a long tail:
$s_X(z)\le L_0$ on its support, while $b(x)\ne0$ forces
$s_X(x)>2L_0$, hence $|s|\ge L_0$.

For $V_\beta$, the derivative support has $|s|a\ge c_\alpha$.
Polynomial factors in $|s|$ therefore give a normal integral at most
$C_\alpha e^{-c_\alpha/a}$.  Freeze $b(x)^2$ at $b(z)^2$ as above:
the frozen contribution is
$C_\alpha e^{-c_\alpha/a}f^2b(z)^2\chi(z)^2$, absorbed by the sheet
gain since $e^{-c_\alpha/a}\le o_\delta(1)a^2\le o_\delta(1)D$.
The freezing error on $S_{\rm sw}$ is at most $C L_0^{-1}D_{\rm sw}$;
off it, the error is a long tail bounded by $o_{L_0}(1)N_R$.

\emph{Cross terms with the partition derivative.}
Only the three displayed crosses against $V_b$ remain.  If
$V_b\cdot V_f\ne0$, then $f'(d(z))\ne0$ and $x$ is in the physical
switch.  The collar plateau gives $s_X(z)>10L_0$, so $|s|>7L_0$.
Young's inequality pays $\varepsilon|V_f|^2$ from the exact derivative
energy and leaves $C_\varepsilon|V_b|^2$ on this long-ray support.
It is a tail controlled by $f\le1$, with no bound on $f'$ required.
The cross $V_b\cdot V_\mu$ is a long tail for the same support reason
as the $V_\mu$ square.  Finally, on $S_{\rm sw}$ the absolute
$V_b\cdot V_\beta$ integral is at most
$C L_0^{-1}\int_{S_{\rm sw}}\chi^2e^{-c_\alpha/a}
\le C L_0^{-1}D_{\rm sw}$; off that set it is again a long tail.
All other crosses, including those between $V_f$ and $V_\beta$, were
included in the first Young bound.

Consequently the costs apart from \eqref{eq:M4} are at most
\[
 (1+\tau')\sigma\int_{\Sigma_g}\chi^2f'^2|\nabla d|^2
 +C_{\tau'}R^{-2}
       \left(\int_{\Sigma_g}\chi^2f^2+C_{L_0}N_R\right)
 +o_{L_0}(1)(D_{\rm sw}+N_R),
\]
plus arbitrarily small fractions of the favorable mass already
retained.  Every occurrence of $f'$ has been assigned to its displayed
derivative energy.  All remaining bounds use $0\le f\le1$, so these
constants are independent even of the finite slope bound
$C_f^{(0)}$, as well as of the later bridge parameters.

\paragraph{Claim 3: cancellation of the partition costs.}
The preceding estimates have assigned every occurrence of $f'$ to its
exact derivative energy.  We must now cancel the two signed costs in
\eqref{eq:M4}, rather than spend a fixed amount of the surface mass on
them.  The required estimate is
\begin{equation}\label{eq:M5}
 \int\bigl(\hbox{the two densities in \eqref{eq:M4}}\bigr)
        -\int\mathcal V
 \le o_{L_0}(1)(D_{\rm sw}+N_R).
\end{equation}
The first integral includes the lifted measure for the second density.
The error is uniform for $R$ sufficiently large and for the entire
class of $f$ in the lemma.  We supply the details, including the order
of the logarithmic cutoffs.
We will cancel the ambient partition cost against one lifted cost at
each physical point, so the choice of that lifted preimage must be
unique.  The remaining short-ray costs will be bounded by defect mass,
and long rays by the bad-net mass.

\emph{Short rays and uniqueness.}
Let $\mathcal W=\{2L_0<s_X<3L_0\}$, choose
$M=C_1\log L_0$, and temporarily fix a small $\xi>0$.  A preimage
$(z,s)$ of $x\in\mathcal W$, $|s|\le M$, is called accurate if
\[
 B_z(s):=D(z)e^{(2\sqrt2+\xi)|s|}\le b_\xi,
\]
where the fixed positive $b_\xi$ will be small.  For $L_0$ large,
$M<L_0/2$, so every such base lies in $S_{\rm sw}$ and $p_0=1$.
On an accurate segment the diffuse estimate with exponential loss
$\xi/2$, and $Tp=-\nabla D/2$, imply
\begin{equation}\label{eq:growth-accurate-ray}
 |p(x)-\nu(z)|+|U(x)-s|
 \le C_\xi D(z)e^{(2\sqrt2+\xi/2)|s|}
 =C_\xi B_z(s)e^{-\xi|s|/2}.
\end{equation}
In detail, along the segment
$|T\nu|\le CD+C\sqrt D\,|p-\nu|$,
$|p(z)-\nu(z)|\le D(z)$, and
$\int_0^{|s|}\sqrt{D(z+t\nu)}\,\dd t\le C_\xi\sqrt{b_\xi}$.
Gronwall followed by integration proves
\eqref{eq:growth-accurate-ray}.  Also
$|s|a(z)\le C\sqrt{b_\xi}|s|e^{-(\sqrt2+\xi/2)|s|}$, so
$\beta=1$ if $b_\xi$ is sufficiently small.

There is at most one accurate preimage of a physical point.  Suppose
$(z,s)$ and $(\widetilde z,\widetilde s)$ are both accurate.  Comparing
each signed depth with $U(x)$ first gives
$|s-\widetilde s|\le C_\xi b_\xi$ and comparability of their residual
exponentials.  Comparing their normals with $p(x)$ and using
$z=x-s\nu(z)$ next gives
\[
 |z-\widetilde z|
 \le C_\xi b_\xi(1+|s|)e^{-\xi|s|/2}\le C'_\xi b_\xi.
\]
Choose $b_\xi$ so that these bases belong to one uniform small graph.
On it $|A|\le C\sqrt{D(z)}$ by local defect comparability.  Tangential
projection of the equality of the two normal images yields
\[
 |\operatorname{proj}_{T_z\Sigma}(\widetilde z-z)|
 \le C|\widetilde s|\sqrt{D(z)}\,|\widetilde z-z|.
\]
The left side is at least $.9|\widetilde z-z|$, whereas the coefficient
on the right is at most $C_\xi\sqrt{b_\xi}<.9$.
Thus the bases and depths coincide.  This proof uses signed depths and
oriented normals and therefore covers preimages from opposite sides.

\emph{Cancellation along the unique accurate ray.}
At an accurate pair, $\nabla s=\nu(z)$ in the inverse normal chart, and
\eqref{eq:growth-accurate-ray} gives
\[
 |q/G-1|+|\nabla\log(G/q)|\le C_\xi B_z.
\]
Retaining only the ambient function and this lifted function in
\eqref{eq:growth-variance}, the variance is at least
\[
 \chi^2\frac{q^2G^2}{\alpha_0^2q^2+b^2G^2}
 \left|\nabla\theta+\alpha_0b\nabla\log(G/q)\right|^2.
\]
The two pure partition squares, in physical measure, equal
$\chi^2(q^2b^2+G^2\alpha_0^2)|\nabla\theta|^2$.
For $A=q^2$, $B=G^2$, their coefficient minus the leading variance
coefficient is exactly
\[
 Ab^2+B\alpha_0^2-\frac{AB}{A\alpha_0^2+Bb^2}
 =\frac{\alpha_0^2b^2(A-B)^2}{A\alpha_0^2+Bb^2}.
\]
Expanding the variance and dropping its remaining negative square,
then combining the two $\nabla\chi$ crosses, bounds the uncancelled
cost by
\begin{equation}\label{eq:growth-accurate-error}
 C_\xi\chi^2G^2B_z
       \bigl(L_0^{-1}+L_0^{-2}+(L_0R)^{-1}\bigr).
\end{equation}
There is no multiplicity factor for this cancellation, by the accurate
pair uniqueness just proved.  Other preimages can only increase the
full variance.

\emph{The remaining rays.}
The cancellation has used at most one lifted branch at each physical
point.  We now bound every other partition cost by defect mass or by
the bad-net mass.  For inaccurate lifted preimages of depth at most $M$,
\[
 G(s)^2\le C_\xi D(z)e^{\xi M}.
\]
If an ambient point has no accurate preimage, select a closest zero.
When its distance is at most $M$, that zero is regular and belongs to
$S_{\rm sw}$, since $X$ is more than $2L_0$ away.  Its ray is normal
and inaccurate, whether or not it belongs to the curvature cutoff.
Indeed an accurate full normal ray has $\beta=1$ by the bound on
$|s|a(z)$ proved above.  Its base in $S_{\rm sw}$ has $\mu=f=1$, so
it would supply an accurate lifted preimage, contrary to the case
under consideration.
The phase estimate \eqref{eq:phase-decay}, with exponent loss at most
$\xi/2$, gives
\[
 q(x)^2\le C_\xi D(z)e^{2\xi M}.
\]
This ambient integral is bounded using the area formula for the full
normal map on $|s|\le M$.  One may restrict it to any measurable
closest-point selection; the resulting integral is no larger than the
integral over all eligible rays.  The normal Jacobian is at most
$C(1+M)^m$.  Its critical images have zero volume by the area formula,
so no inverse chart or unique nearest point is required for this
estimate.  Measurable selection exists since the zero set is closed
and nearest points form a nonempty compact set locally.

If the closest distance exceeds $M$, use phase decay and the volume
bound in \eqref{eq:growth-collar-count}.  For these long ambient tails
use one fixed coarse exponent in \eqref{eq:phase-decay}, independently
of $\xi$; only the short-ray comparison above needs an exponent loss
chosen in terms of $\xi$.  For lifted rays of depth
greater than $M$, a contribution in $\mathcal W$ requires
$s_X(z)\le3L_0+|s|$; the same collar bound and the exponential tail of
$G^2$ apply.  Thus, for fixed dimensional $N$ and $c>0$, independent
of sufficiently small $\xi$ and of $C_1$, all costs in
this comparison are bounded by
\begin{equation}\label{eq:growth-log-errors}
 C_\xi\frac{(1+M)^N e^{2\xi M}}{L_0}D_{\rm sw}
 +C_\xi(1+L_0+M)^N e^{-cM}N_R.
\end{equation}
The Jacobian, the integration over depth, and derivatives of $\beta$
only alter the fixed power $N$.  Weight shifts cost $e^{2|s|/R}$ and
are absorbed in the tail exponent for all sufficiently large $R$.

\emph{Fixing the logarithmic window.}
The order of choices is now decisive.  The defect threshold and its
positive bad-mass coefficient $c_\delta$ have already been fixed.
Choose $C_1>(N+3)/c$, then $\xi>0$ with $2\xi C_1<1/4$, then
$b_\xi>0$ as above, and finally let $L_0$ be large.  Both coefficients
in \eqref{eq:growth-log-errors} tend to zero; the first is bounded by a
constant times $L_0^{-3/4}(1+\log L_0)^N$ and the second by a constant
times $L_0^{-3}(1+\log L_0)^N$.  They can therefore be absorbed by the
already fixed positive multiples of $D_{\rm sw}$ and $N_R$.
This proves \eqref{eq:M5} with a fixed threshold $\delta$.

\paragraph{Conclusion of the gluing estimate.}
Claim 1 supplies the positive reserve, Claim 2 leaves only the signed
partition costs, and Claim 3 absorbs those costs with the variance.
Stability of $F$ and division by $\sigma$ therefore give \eqref{eq:M3}.  The remaining $C_{L_0}R^{-2}N_R$ is
absorbed for $R$ large; the interface tolerances and Young constants
were chosen with room inside $\tau$.  This also proves the asserted
uniformity in the center and in $f$.
\end{proof}

\subsection{A reach multiplier with a long plateau}

The letter $P_*$ below denotes the multiplier, and is unrelated to the
fixed integration depth $P$ defining $I_P$.
We now choose $f$ so that its derivative energy in \eqref{eq:M3} can be
estimated by integrating the reach inequality.  The plateau keeps this
derivative away from the collar; a slow initial decay makes the
multiplier flux there arbitrarily small.  The subsequent power and
exponential ranges supply the two different bounds needed for short
and long reaches.

\begin{lemma}[A weight and its reach multiplier]
\label{lem:growth-multiplier}
Let $\zeta_0=.01$ and $k=1.53$.  Increase $t_{\min}$, before fixing
$\delta$, so that
$\exp(C/(\gamma t_{\min}^{\gamma}))<1.001$, where $C,\gamma$ dominate
the errors in \eqref{eq:C8}.  For every finite $t_*\ge t_{\min}$,
every $\eps_3>0$, and every $\eta>0$, there are $t_0>t_*$ and a
function $f$ as in Lemma~\ref{lem:growth-gluing} such that, with
\[
 B_*(t)=1+Ct^{-\gamma}+\zeta_0\mathbf1_{\{t\ge t_0\}},\quad
 p_*(t)=\exp\int_{t_{\min}}^t\frac{B_*(s)}s\,\dd s,\quad
 P_*(t)=\frac1{p_*(t)}\int_t^\infty p_*(s)f'(s)^2\,\dd s,
\]
one has $-P_*'=f'^2+B_*P_*/t$ almost everywhere and
\begin{equation}\label{eq:M6}
 \begin{gathered}
 tP_*(t)\le .79f(t)^2\qquad(t_{\min}\le t<t_0),\\
 P_*(t)\le(\kappa/4+.001)f(t)^2,\qquad
 t^3f(t)^2\le\eps_3\qquad(t\ge t_0),\\
 \sup_{[t_{\min},t_*]}P_*\le\eta.
 \end{gathered}
\end{equation}
Moreover
\[
 e^{-\kappa t}\le f(t)^2\le C_f e^{-\kappa t},\qquad
 \Psi(t):=\int_t^\infty P_*(s)\,\dd s\le C_f f(t)^2
\]
for a finite constant $C_f$ depending on the final construction.
\end{lemma}

\begin{proof}
There are four ranges.  A plateau protects the collar; a very slow
power decay makes the multiplier flux on that plateau small; a faster
power decay makes $t^3f^2$ small; and the final exponential tail matches
$E=e^{-\kappa d}$.  Take $t_s\ge t_*$ and use the continuous piecewise
formula
\[
 f(t)=
 \begin{cases}
 1,&t\le t_s,\\
 (t/t_s)^{-\mu},&t_s<t<t_p,\\
 f(t_p)(t/t_p)^{-k},&t_p<t<t_e,\\
 f(t_e)e^{-\kappa(t-t_e)/2},&t\ge t_e.
 \end{cases}
\]
Here $0<\mu\ll1$, $t_p>t_s$, and $t_p<t_0<t_e$.  Corners cause no
problem for locally Lipschitz tests or the multiplier identity.
Set $F_0=\exp(C/(\gamma t_{\min}^{\gamma}))<1.001$.
For $s\ge t$,
\[
 \frac{p_*(s)}{p_*(t)}
 \le F_0(s/t)^{1+\zeta_0}.
\]
\paragraph{The main power range and the exponential tail.}
On the main polynomial bridge its polynomial contribution is therefore
at most $F_0 A/t$ in $P_*/f^2$, where
\[
 A=\frac{k^2}{2k-\zeta_0},\qquad F_0A<.77.
\]
The tail value $Q_e=P_*(t_e)/f(t_e)^2$ obeys
\[
 Q_e\le
 \frac{\kappa^2/4}
 {\kappa-(1+\zeta_0+Ct_e^{-\gamma})/t_e}
 =\kappa/4+o_{t_e\to\infty}(1).
\]
The same upper bound holds for $P_*(t)/f(t)^2$ for $t\ge t_e$.
Propagating this terminal flux backwards along the polynomial range
adds at most
\[
 F_0Q_e(t_e/t)^{1+\zeta_0-2k}
\]
to $P_*/f^2$.  Its product with $t$ tends uniformly to zero on
$[t_p,t_0]$ as $t_e\to\infty$ with $t_0$ fixed, since
$1+\zeta_0-2k=-2.05$.
For $t_0\le t\le t_e$, the total is at most
$F_0(A/t_0+Q_e)$.  This is less than $\kappa/4+.001$ by first taking
$t_0$ large and then $t_e$ large: the fixed factor costs less than
$.000354$ since $\kappa\le\sqrt2$, leaving a positive margin for the
two vanishing errors.  This proves both multiplier bounds on the main
bridge and the tail.

\paragraph{Small flux on the collar plateau.}
On the slow bridge there is no $\zeta_0$ drift.  Up to the factor
$F_0$, the ratio $tP_*/f^2$ is bounded by
\[
 \frac\mu2\bigl(1-(t/t_p)^{2\mu}\bigr)
 +(t/t_p)^{2\mu}\frac{t_pP_*(t_p)}{f(t_p)^2}.
\]
Thus it remains below $.79$, and its value at $t_s$ can be made as
small as desired by taking $\mu$ small and then $t_p/t_s$ large.
On the plateau $tP_*(t)\le F_0 t_sP_*(t_s)$, proving the arbitrarily
small compact flux.  More explicitly, require the terminal contribution
to $tP_*/f^2$ on $[t_p,t_0]$ to be at most $.005$.  The main, slow,
and plateau ranges then have the common upper bound
$F_0^2(F_0A+.005)<.775<.79$, and
\[
 \sup_{[t_{\min},t_*]}P_*
 \le \frac{F_0^2}{t_{\min}}
       \left(\frac\mu2+.775(t_s/t_p)^{2\mu}\right).
\]
All these choices can be made before selecting
$t_0,t_e$, since the later terminal contribution can be required to be
smaller than any prescribed positive error on $[t_{\min},t_p]$.

\paragraph{Completing the choices.}
The strict inequality $2k>3$ supplies the last decay requirement:
after $t_p$, $t^3f(t)^2$ decreases as $t^{3-2k}$.  Choose $t_0$ large
enough to put it below $\eps_3$, in addition to the preceding lower
bounds.  Choose $t_e$ afterwards; on its exponential tail the same
quantity is decreasing once $t_e>3/\kappa$.  All the required bounds
are simultaneous.  Finally $-f'/f\le k/t<\kappa/2$ before $t_e$,
after increasing $t_{\min}$.  Hence $f^2e^{\kappa t}$ is nondecreasing
up to $t_e$ and constant thereafter, proving both exponential
comparisons.  The primitive estimate follows on the tail by integration
and on the remaining compact interval by positivity of $f$.
\end{proof}

\subsection{Integrating the reach inequality and paying for the feet}

There are two kinds of cost to compare with the surface mass in
\eqref{eq:M3}: a source at the point where the reach is differentiated,
and a positive charge at a receiving foot of a plateau contact.
We estimate the local sources first, then integrate the inequality and
transfer the foot charges using their multiplicity-one bound.

\paragraph{Local source costs.}
Choose the constants in \eqref{eq:C8} so that $B_*$ dominates its drift.
On nonplateau branches with $d<t_0$, \eqref{eq:M6} bounds the multiplied
curvature source by $.80 f^2a^2$, after increasing $t_{\min}$ and making
the already fixed reach tolerance small.  On $d\ge t_0$, if
$ad\le c_2$ for a sufficiently small fixed $c_2>0$, the gradient lower
bound in \eqref{eq:C8} gives
\[
 (1+\eta_0+Cd^{-\gamma})a^2d
 \le\zeta_0\frac{|\nabla d|^2}{d}.
\]
This is paid for by the additional drift in $B_*$.  If $ad>c_2$,
the multiplied source is instead at most
\[
 C f^2a^2d\le C\eps_3 a^2d^{-2}
 \le C_{c_2}\eps_3 a^4.
\]
Plateau branches have no such curvature source.  Their outgoing source
in \eqref{eq:C9} satisfies
\[
 G_1P_*|H|\le .13 f^2 I.
\]
On the early range use $P_*/f^2\le .79/t_{\min}$; on the late range
use \eqref{eq:M6} and \eqref{eq:growth-sharp-budget}.  The limiting
late coefficient is
$(\kappa/4)G_1/(2\sqrt2)=1/8$, strictly below $.13$.

\paragraph{The weak integration and its cutoff errors.}
Let $m_0$ be an exterior Lipschitz cutoff, zero for $s_X\le5L_0$ and
one for $s_X\ge6L_0$, or identically one if $X$ is empty.  Multiply the
measure reach inequality by $P_*(d)m_0\chi^2$.  For its absolutely
continuous part, write $\Delta d\le A+B|\nabla d|^2/d$.  Integration
by parts and $-P_*'=f'^2+B_*P_*/d$ give
\begin{equation}\label{eq:growth-multiplier-integration}
 \begin{split}
 &\int m_0\chi^2 f'^2|\nabla d|^2
 +\int m_0\chi^2(B_*-B)P_*\frac{|\nabla d|^2}{d}\\
 &\hspace{12mm}\le\int AP_*m_0\chi^2
       +\int P_*\nabla d\cdot\nabla(m_0\chi^2).
 \end{split}
\end{equation}
The nonpositive singular part of $\Delta d$ has the favorable sign.
Nonnegative compactly supported Lipschitz tests are valid by uniform
approximation and convergence of their weak gradients; locally $d$ is
semiconcave and $\nabla d$ is bounded.  Chain rules apply almost
everywhere across the finitely many coefficient corners.  Since
$f'=0$ on the whole collar $s_X\le10L_0$, the first integral in
\eqref{eq:growth-multiplier-integration} is precisely the gradient
integral in \eqref{eq:M3}.

The cutoff derivative in \eqref{eq:growth-multiplier-integration} is
bounded as follows.  The part with $\nabla m_0$ costs
\[
 C_{\delta,L_0}\sup_{[t_{\min},t_*]}P_*\,N_R
\]
by
\eqref{eq:C7} and \eqref{eq:growth-bounded-collar-reach}.  On the other
part use $\nabla\Psi(d)=-P_*(d)\nabla d$ and integrate once more:
\[
 \int m_0P_*\nabla d\cdot\nabla\chi^2
 =\int\Psi(d)\bigl(m_0\Delta_\Sigma\chi^2
                      +\nabla m_0\cdot\nabla\chi^2\bigr).
\]
For $r_c=\sqrt{1+|x-x_c|^2}$, convexity of $r_c$ gives the one-sided
ambient bound $\nabla^2\chi^2\le CR^{-2}\chi^2\Id$.  Its tangential
trace and $\Delta_\Sigma x=-H\nu$ therefore give
\[
 \Delta_\Sigma\chi^2\le CR^{-2}\chi^2+CR^{-1}D\chi^2.
\]
Together with the primitive bound, the cutoff derivative is at most
\begin{equation}\label{eq:growth-spatial-error}
 \begin{split}
 &C_{\delta,L_0}\sup_{[t_{\min},t_*]}P_*\,N_R
 +C_fR^{-2}\int_{\Sigma_g}\chi^2f^2\\
 &\quad+C_fR^{-1}\int_{\{s_X\ge4L_0\}}\chi^2Df^2
 +C_{\delta,L_0,f}R^{-1}N_R.
 \end{split}
\end{equation}
Exhaustion of these identities is justified by the same exponential
estimates as in the gluing lemma.  In particular no sign of a second
derivative of $m_0$ is required.

\paragraph{Incoming positive charges.}
The local sources and the cutoff errors have now been accounted for.
It remains to pay the mean-curvature term at a second contact point.
We account for the positive foot charge of every selected plateau
pair $(z,y)$.  Choose a measurable active branch, selecting a plateau
whenever one is available.  First the part $d(z)>\sqrt R$ costs
$o(K(R)R^{m-4})$ directly at its source.  Indeed $P_*$ is eventually
bounded by $C_fe^{-\kappa d}$; the sources in
\eqref{eq:C8}--\eqref{eq:C9} and the bound \eqref{eq:C7} grow at most
polynomially in $d$, while the tangent-ball bound controls the foot
mean curvature by $C/d$.  The good weighted area is bounded by
\eqref{eq:growth-area}.  Thus this error is at most
$C_fe^{-c\sqrt R}K(R)R^m$ after altering $c>0$.

On the remaining pairs $|z-y|\le d(z)\le\sqrt R$, so
$\chi(z)^2/\chi(y)^2\le e^{2/\sqrt R}=1+o(1)$ uniformly in the
center.  If $s_X(y)\le10L_0$, the source reach is at most $t_*$ by
\eqref{eq:growth-bounded-collar-reach}, and $z$ is in a fixed bounded
enlargement of $X$.  Its entire cost is
\[
 C_{\delta,L_0,t_*}
       \sup_{[t_{\min},t_*]}P_*\,N_R,
\]
where the coefficient is fixed before $f$ is constructed.
For every other receiving foot apply \eqref{eq:C10}.  Its selected
positive charge has inverse multiplicity at most one and area factor
\[
 J_{\rm foot}=\frac{|\det(1+\ell A_{N_y}(y))|}{s}\le C_\Gamma,
 \qquad d(y)\le d(z),\qquad \ell=d(z)/G_1.
\]
On $d(z)<t_0$, its cost is at most $.14\chi(y)^2 f(d(y))^2I(y)$
by choosing $t_{\min}$ large after the fixed bound $C_\Gamma$.
On $d(z)\ge t_0$, if $a(y)d(z)\le c_3$ with sufficiently small fixed
$c_3$, the improved area factor is less than $1.05$.  Equations
\eqref{eq:M6} and \eqref{eq:growth-sharp-budget}, monotonicity of $f$,
and the weight ratio then give the same $.14$ bound.  Its limiting
coefficient is $1.05/8=.13125$, leaving strict room.  If instead
$a(y)d(z)>c_3$, use $H_{N_y}(y)_+\le C/d(z)$ and $d(z)^3f(d(z))^2
\le\eps_3$ to bound the transferred cost by
\[
 C_{\Gamma,c_3}\eps_3\chi(y)^2a(y)^4.
\]
Every foot outside the bounded-collar case lies in $s_X>10L_0$, hence
in the region of full strength in \eqref{eq:M3}.  The count of positive
charges, not a two-sided unoriented count, is what permits the single
coefficient $.14$.

\paragraph{The numerical reserve.}
At a given surface point the selected reach branch contributes either
a nonplateau source or a plateau source.  It may also receive a positive
charge, but the transfer counts that charge at most once.  After
factoring $\chi^2f^2$, the combined costs are
\[
 \begin{array}{c|c}
 \text{selected local branch}&\text{local and incoming costs}\\ \hline
 \text{nonplateau}& .80a^2+.14I\\
 \text{plateau}& .27I.
 \end{array}
\]
A point receiving only a charge has a smaller cost.  For
$\tau\le .001$, \eqref{eq:growth-sharp-budget} implies
\begin{equation}\label{eq:growth-numerical-margin}
 \begin{aligned}
 (1-\tau)\frac{a^2+I}{2}
       -(1+\tau)(.80a^2+.14I)&\ge .05I,\\
 (1-\tau)\frac{a^2+I}{2}-(1+\tau).27I&\ge .22I.
 \end{aligned}
\end{equation}
For example the first left side is at least
$[.35936-.3013/.999]I>.057I$ at the worst allowed tolerance.
Thus the sharp comparison $I\ge .999a^2$ is essential here.

Combine \eqref{eq:M3} and
\eqref{eq:growth-multiplier-integration}, use
\eqref{eq:growth-numerical-margin} and $I\ge c_ID$, and absorb the
two $R^{-1}$ terms in \eqref{eq:growth-spatial-error} after $f$ is fixed.
Choose the compact multiplier flux so small that all fixed collar
costs are absorbed by $c_0N_R$.  This proves
\begin{equation}\label{eq:M7}
 c_4\left(N_R+
      \int_{\{s_X\ge4L_0\}}\chi^2Df^2\right)
 \le\eps_4 A_4+C_fR^{-2}\int_{\Sigma_g}\chi^2f^2
       +C_fK(R)R^{m-4}.
\end{equation}
For example, after taking the compact flux small enough to spend at
most $c_0N_R/4$ and then absorbing the $R^{-1}$ terms, one may use
$c_4=\min\{c_0/4,.02c_I\}>0$.  This constant is fixed by the gluing
mass and numerical margins and is independent of $f$.
The coefficient $\eps_4$ is any prescribed
positive number: choose $\eps_3$ in Lemma~\ref{lem:growth-multiplier}
after fixing $c_4$ and the preceding source constants.  The terms
denoted $C_f$ may then be large, but remain fixed as $R$ grows.

\subsection{Closure of the estimate}

We have estimated curvature by a small multiple of the squared defect,
and we have estimated the first defect moment weighted by $f^2$.
The remaining step is to connect these estimates.  The integral
comparison \eqref{eq:L11} reduces squared defect and interaction to
curvature and $DE$.  Since $E\le f^2$, the latter term is paid by
\eqref{eq:M7}.  We perform this absorption before allowing the final
comparison constant $C_f$ to enter the inverse-square cutoff error.

Set $Z_2=\int_{\Sigma_g}\chi^2(D^2+E^2)$ and
$V_R=K(R)R^{m-4}$.  Equations~\eqref{eq:L11} and \eqref{eq:M2}, with
$\eps_2 C_D$ small, imply
\begin{equation}\label{eq:growth-preclosure}
 A_4+Z_2+N_R
 \le C'\left(N_R+\int_{\Sigma_g}\chi^2DE+V_R\right).
\end{equation}
The constant $C'$ is fixed before constructing $f$.  Since $E\le f^2$
and the good collar has weighted area at most $C_{L_0}N_R$,
\[
 \int_{\Sigma_g}\chi^2DE
 \le\int_{\{s_X\ge4L_0\}}\chi^2Df^2+C_{L_0}N_R.
\]
Choose $\eps_4$ in \eqref{eq:M7} small in terms of $C',c_4,C_{L_0}$.
Substitution in \eqref{eq:growth-preclosure} absorbs the resulting
multiple of $A_4$.  The remaining inverse-square term satisfies,
for any $\eta>0$,
\[
 C_fR^{-2}\int_{\Sigma_g}\chi^2f^2
 \le C_f'R^{-2}\int_{\Sigma_g}\chi^2E
 \le\eta\int_{\Sigma_g}\chi^2E^2
       +C_{f,\eta}R^{-4}\int_{\Sigma_g}\chi^2
 \le\eta Z_2+C_{f,\eta}V_R.
\]
Take $\eta$ small enough for the last absorption.  This proves
\eqref{eq:M1} and Proposition~\ref{prop:growth-fourth-order}.
Every shift of a weight above was controlled by the Lipschitz bound
$|\nabla\log\chi|\le R^{-1}$; consequently none of the constants
depends on the translating center $x_c$.
